\problema{Merge k Sorted Lists}{23}{src/04-heaps-topk/23-merge-k-sorted-lists.cpp}{H}

Nos dan un arreglo de \texttt{k} listas enlazadas, cada una ya ordenada de
menor a mayor. Hay que fusionarlas en una sola lista enlazada, también
ordenada, y devolver su cabeza.

\paragraph{La idea, con un ejemplo de la vida real.} Fusionar dos listas
ordenadas ya es un ejercicio clásico: comparas las dos cabezas, te quedas
con la más chica, avanzas esa lista, y repites. Aquí el reto es que hay
\texttt{k} listas en vez de dos. Imagina \texttt{k} filas de personas, cada
fila ya ordenada de menor a mayor estatura, y quieres formar una sola fila
combinada ordenada. En cada paso necesitas saber cuál de las \texttt{k}
personas al frente de cada fila es la más bajita; compararlas una por una
cada vez que avanzas sería lento si \texttt{k} es grande. En vez de eso,
metes a las \texttt{k} personas que están al frente de cada fila en un
\emph{min-heap}: sacar a la más bajita de un grupo de \texttt{k} candidatos
cuesta mucho menos que compararlas todas cada vez.

\begin{center}
\ttfamily
\begin{tabular}{l}
lists = [[1,4,5],[1,3,4],[2,6]]\\
fusión: 1,1,2,3,4,4,5,6\\
\end{tabular}
\end{center}

\paragraph{Cómo lo resuelve el código, paso a paso.}
\begin{enumerate}
  \item Metemos al min-heap la cabeza de cada lista no vacía (el primer nodo
        de cada una).
  \item Creamos un nodo \texttt{dummy} (ancla vacía) para ir construyendo la
        lista de respuesta, y un puntero \texttt{cola} que empieza
        apuntando a ese \texttt{dummy}.
  \item Mientras el heap no esté vacío: sacamos el nodo con menor valor,
        lo enganchamos al final de la respuesta (\texttt{cola->next}) y
        avanzamos \texttt{cola}.
  \item Si el nodo que acabamos de sacar tenía un siguiente (su lista
        original no se había terminado), metemos ese siguiente al heap: se
        vuelve el nuevo candidato de esa lista.
  \item Devolvemos \texttt{dummy->next}: el primer nodo real de la fusión.
\end{enumerate}
En cualquier momento, el heap contiene exactamente un candidato por cada
lista que todavía tiene nodos sin usar (su nodo actual, el más pequeño no
usado de esa lista en particular, porque las listas ya vienen ordenadas). El
menor de todos esos candidatos es, por lo tanto, el menor valor no usado de
\emph{todas} las listas juntas, así que siempre es seguro colocarlo a
continuación.

\lstinputlisting{src/04-heaps-topk/23-merge-k-sorted-lists.cpp}

\complejidad{$O(N \log k)$}{$O(k)$}

\paragraph{¿Qué significa esa complejidad?} $N$ es el número total de nodos
en todas las listas juntas, y \texttt{k} es el número de listas. Cada uno de
los $N$ nodos entra y sale del heap exactamente una vez, y cada una de esas
operaciones cuesta $\log k$ porque el heap nunca tiene más de \texttt{k}
elementos a la vez (uno por lista activa). El espacio extra es solo
$O(k)$ para el heap, ya que reutilizamos los nodos originales para construir
la lista de salida en vez de crear nodos nuevos.

\paragraph{Consejos para la entrevista (Amazon).} Este problema combina dos
técnicas que suelen evaluar por separado: fusión de listas enlazadas y uso
de heaps para elegir entre múltiples candidatos. Vale la pena mencionar la
alternativa de fusionar las listas de dos en dos (fusión por pares, similar
a merge sort), que también logra $O(N \log k)$ pero sin necesitar un heap
explícito. Un detalle de implementación en C++ que suelen preguntar: cómo
comparar punteros a nodos dentro de un \texttt{priority\_queue} usando un
comparador personalizado, ya que no se puede comparar directamente
\texttt{ListNode*} por valor. También aclara qué pasa si alguna lista viene
vacía (\texttt{nullptr}) o si el arreglo completo de listas viene vacío: en
ambos casos la respuesta debe ser \texttt{nullptr}, sin caer en un acceso
inválido.
